% Experimental results.
\providecommand{\tbd}{\textemdash}

\section{Evaluation}
\label{sec:evaluation}

\subsection{Regularization comparison on MNIST with a fully connected network}
\label{subsec:reg-mnist-fcn}
Table~\ref{tab:regularization-mnist-fcn} compares regularization methods after selecting each method's own hyperparameters on a \emph{validation} split using a unified protocol. Everything below is shared across methods except the meaning of each method's grid point (the knob being swept).

\paragraph{Validation-based hyperparameter protocol (GP Bayesian optimization).}
Hyperparameters are selected on a validation split using Gaussian-Process (GP) Bayesian optimization; the test set is used only for final numbers in Table~\ref{tab:regularization-mnist-fcn}. Here, Bayesian optimization (BO) is a sample-efficient search method for expensive black-box objectives (mean validation accuracy after training): it uses a GP surrogate to predict promising regions and an Expected-Improvement acquisition rule to balance exploration and exploitation, and one BO step in this report means evaluating one candidate configuration across the fixed seed list. The protocol is:
\begin{itemize}
  \item \textbf{Shared training protocol:} identical dataset splits, architecture, optimizer type, learning rate, batch size, maximum epochs, and early-stopping rule (when applicable), unless a row is explicitly ablating one of these.
  \item \textbf{Validation-only selection:} we hold out a validation subset of the training data (e.g.\ $\ValSplitPct\%$ of training examples, stratified by class). Bayesian optimization targets mean validation accuracy; test labels are never used during search.
  \item \textbf{Globally fixed random seeds (not tuned):} before any search, we fix one seed list (here: $\{\SeedList\}$, i.e., \NumSeeds{} seeds). Every BO candidate for every method is trained once per seed in this list. Seeds are not hyperparameters and are never selected for performance.
  \item \textbf{Surrogate model and acquisition:} for each method, we fit a GP surrogate over that method's hyperparameter space (continuous knobs are optimized in log space where appropriate) and select the next trial by Expected Improvement.
  \item \textbf{Equal BO trial budget per tunable method:} each tunable method receives the same number of BO candidates (\BOBudgetTotal{} total: \BOBudgetInit{} initial random points + \BOBudgetGuided{} GP-guided points). One candidate evaluated across the fixed seed list is counted as one BO step. The plain baseline has no tunable regularization coefficient and is therefore evaluated once (with the same seed list).
  \item \textbf{Search domains:} $L_2$ decay coefficient $\lambda_2 \in [10^{-5}, 3\cdot10^{-3}]$, $L_1$ penalty coefficient $\lambda_1 \in [10^{-5}, 3\cdot10^{-3}]$, dropout rate $p \in [0,0.5]$, early-stopping patience $K \in [5,20]$ (integer), and balancing schedule treated as a categorical variable over predefined modes (no balancing, pre-balance only, and selected interleaving frequencies).
  \item \textbf{Selection and reporting:} for each method, the chosen configuration is the one with the highest mean validation accuracy across the fixed seed list among all BO-evaluated candidates. The table reports test metrics for that chosen configuration only, with uncertainty (e.g.\ $\sigma_{\mathrm{test}}$) computed across the same seeds.
\end{itemize}

\paragraph{Limitation (BO still has design choices).}
Using GP BO avoids hand-crafted fixed grids, but fairness still depends on choices such as search-domain bounds, acquisition function, initialization strategy, and total BO budget. Different choices can change the selected hyperparameters and, in turn, final ranking.

\paragraph{Fixed experimental setting.}
Dataset: MNIST (official train/test split; $\MnistTrainSize{}$ training and $\MnistTestSize{}$ test images).
Inputs: grayscale $28\!\times\!28$ flattened to $784$; ToTensor scaling to $[0,1]$ followed by per-channel normalization to $[-1,1]$ (mean $0.5$, std.\ $0.5$), as implemented in the code.
Architecture: \texttt{small\_fcn} fully connected network, $784 \rightarrow 256 \rightarrow 10$, ReLU activations on the hidden layer, linear output (10 classes).
Optimization: stochastic gradient descent (SGD) without momentum, learning rate $\LearningRate$, batch size $\BatchSize$.
Regularization control: each row uses the backbone above with regularizer hyperparameters \emph{selected on the validation split} by GP BO in this subsection (plain uses $\lambda_2=0$, no dropout, no balancing).
Training: maximum $\MaxEpochs$ epochs per run; rows without early stopping train for the full budget; the early-stopping row uses the patience $K$ selected on validation by BO.
Target for convergence speed: test accuracy threshold $\tau = \TargetTauPct\%$; \emph{Epochs@$\tau$} is the first epoch (1-based index) at which test accuracy reaches or exceeds $\tau$.
Reproducibility: the seed list is fixed \emph{a priori} (see above); $\sigma_{\mathrm{test}}$ is the standard deviation of final test accuracy across those seeds for the selected configuration.

% Controlled-setting comparison: replace \tbd with numbers when available.
\providecommand{\tbd}{\textemdash}

\begin{table}[t]
  \centering
  \small
  \setlength{\tabcolsep}{4pt}
  \caption{%
    Regularization methods under one controlled MNIST + fully connected setting (details in Section~\ref{subsec:reg-mnist-fcn}).
    Rows vary only the regularization strategy; all other protocol choices match.
    For BO-selected rows, scalar entries report means across the fixed seed list; the plain row reports the available single run.
    \emph{Epochs@$\tau$}: first epoch (1-based) at which test accuracy reaches target $\tau$ (here $\tau=\TargetTauPct\%$; see Section~\ref{subsec:reg-mnist-fcn}).
    \emph{Best val}: best validation accuracy observed during training.
    \emph{Gap}: final train accuracy minus final test accuracy (percentage points).
    \emph{$\sigma_{\mathrm{test}}$}: standard deviation of final test accuracy across random seeds.%
  }
  \label{tab:regularization-mnist-fcn}
  \begin{tabular}{@{}lcccccccc@{}}
    \toprule
    Method
    & Epochs@$\tau$
    & \begin{tabular}{@{}c@{}}Final\\train acc.\end{tabular}
    & \begin{tabular}{@{}c@{}}Final\\train loss\end{tabular}
    & \begin{tabular}{@{}c@{}}Final\\test acc.\end{tabular}
    & \begin{tabular}{@{}c@{}}Final\\test loss\end{tabular}
    & \begin{tabular}{@{}c@{}}Best\\val acc.\end{tabular}
    & Gap
    & $\sigma_{\mathrm{test}}$ \\
    \midrule
    Plain (no explicit reg.)
    & \RegMnistPlainEpochsAtTau & \RegMnistPlainFinalTrainAccPct & \RegMnistPlainFinalTrainLoss & \RegMnistPlainFinalTestAccPct & \RegMnistPlainFinalTestLoss & \RegMnistPlainBestValAccPct & \RegMnistPlainGapPpts & \RegMnistPlainSigmaFinalTestAccPct \\
    $L_2$ weight decay
    & \RegMnistLTwoEpochsAtTau & \RegMnistLTwoFinalTrainAccPct & \RegMnistLTwoFinalTrainLoss & \RegMnistLTwoFinalTestAccPct & \RegMnistLTwoFinalTestLoss & \RegMnistLTwoBestValAccPct & \RegMnistLTwoGapPpts & \RegMnistLTwoSigmaFinalTestAccPct \\
    $L_1$ penalty
    & \RegMnistLOneEpochsAtTau & \RegMnistLOneFinalTrainAccPct & \RegMnistLOneFinalTrainLoss & \RegMnistLOneFinalTestAccPct & \RegMnistLOneFinalTestLoss & \RegMnistLOneBestValAccPct & \RegMnistLOneGapPpts & \RegMnistLOneSigmaFinalTestAccPct \\
    Dropout
    & \RegMnistDropoutEpochsAtTau & \RegMnistDropoutFinalTrainAccPct & \RegMnistDropoutFinalTrainLoss & \RegMnistDropoutFinalTestAccPct & \RegMnistDropoutFinalTestLoss & \RegMnistDropoutBestValAccPct & \RegMnistDropoutGapPpts & \RegMnistDropoutSigmaFinalTestAccPct \\
    Early stopping (same budget)
    & \RegMnistEarlyStopEpochsAtTau & \RegMnistEarlyStopFinalTrainAccPct & \RegMnistEarlyStopFinalTrainLoss & \RegMnistEarlyStopFinalTestAccPct & \RegMnistEarlyStopFinalTestLoss & \RegMnistEarlyStopBestValAccPct & \RegMnistEarlyStopGapPpts & \RegMnistEarlyStopSigmaFinalTestAccPct \\
    Synaptic balance ($L_1$)
    & \RegMnistSynBalLOneEpochsAtTau & \RegMnistSynBalLOneFinalTrainAccPct & \RegMnistSynBalLOneFinalTrainLoss & \RegMnistSynBalLOneFinalTestAccPct & \RegMnistSynBalLOneFinalTestLoss & \RegMnistSynBalLOneBestValAccPct & \RegMnistSynBalLOneGapPpts & \RegMnistSynBalLOneSigmaFinalTestAccPct \\
    Synaptic balance ($L_2$)
    & \RegMnistSynBalLTwoEpochsAtTau & \RegMnistSynBalLTwoFinalTrainAccPct & \RegMnistSynBalLTwoFinalTrainLoss & \RegMnistSynBalLTwoFinalTestAccPct & \RegMnistSynBalLTwoFinalTestLoss & \RegMnistSynBalLTwoBestValAccPct & \RegMnistSynBalLTwoGapPpts & \RegMnistSynBalLTwoSigmaFinalTestAccPct \\
    \bottomrule
  \end{tabular}
\end{table}


\paragraph{Main takeaways.}
Among the BO-selected methods, synaptic balance ($L_1$) is the clear standout in this controlled setting: it reaches $\tau$ in \RegMnistSynBalLOneEpochsAtTau epochs and finishes at \RegMnistSynBalLOneFinalTestAccPct\% final test accuracy, which is \RegMnistSynBalLOneGainOverBestNonBalancePpts percentage points above the strongest non-balance baseline (\RegMnistBestNonBalanceFinalTestAccPct\%). Synaptic balance ($L_2$) is also the fastest among the non-$L_1$ alternatives and still improves on the strongest non-balance baseline by \RegMnistSynBalLTwoGainOverBestNonBalancePpts percentage points.

The conventional regularizers cluster much more tightly: $L_1$ penalty, $L_2$ weight decay, and early stopping all finish around 91.96\% test accuracy, while dropout is the best non-balance baseline at \RegMnistDropoutFinalTestAccPct\%. Their train--test gaps are slightly negative, suggesting mild underfitting in this setting, whereas synaptic balance ($L_1$) keeps the best held-out performance despite a positive gap of \RegMnistSynBalLOneGapPpts percentage points. The plain baseline currently has only \RegMnistPlainSeedCount completed seed, so its $\sigma_{\mathrm{test}}$ entry is intentionally left blank.


